\chapter{Introduction}
In the early universe, there were small irregularities in the density of matter. Today, these irregularities can be seen in the Cosmic Microwave Background Radiation \citep{durrer2015}. When gravity starts to act on these small fluctuations, regions that have a slightly higher density will exert a stronger gravitational pull, attracting more matter. This gravitational pull starts off a process where the dense regions will grow denser and the less dense regions will get emptier. Over time, this process leads to the formation of over-dense regions, clusters, and under-dense regions, voids \citep{zeldovich1970}. The formation of these structures have been studied extensively with the use of hydrodynamical simulations \citep{schaye2014, crain2015, schaye2023, kugel2023}, as it is difficult to observe this structure formation directly \citep{bos2016}. \par 
These hydrodynamical simulations have an initial and final particle distribution, often with in-between time-steps. This provides a natural framework for Optimal Transport theory (OT), a rapidly developing area of mathematics. OT allows us to compute distances and relationships between probability distributions \citep{villani2003}. In OT, one looks for the transport of some type of particles (from the initial to the final positions), and tries to minimize the cost this transport takes, making use of some cost function. This cost function is often taken as the work it takes to move particles, possibly with other transport costs, such as fuel costs \citep{villani2008}. \par 
In simulations governed by physical laws, the exact form of cost function is typically unknown, due to the complex gravitational dynamics and non-linear evolution of structures. The reconstruction of initial conditions has been studied within the field of cosmology \citep{nikakhtar2024}, but the inverse OT problem has to our knowledge not been applied in the setting of structure formation in cosmology. Therefore, the goal of this thesis is to reverse the optimal transport process; can we infer the cost function from the observed particle movements?  \par 
To reverse the OT process, data of the FLAMINGO (Full-hydro Large-scale structure simulations with All-sky Mapping for the Interpretation of Next Generation Observations) simulations will be used, see also \citet{schaye2023, kugel2023}. There are different algorithms and methods to recover the cost matrix, each designed for different constraints \citep{stuart2020, ma2021, li2019, chiu2022}. In this thesis, a modified algorithm for unconstrained convex optimization of inverse OT by \citet{ma2021} will be used. \par 
Once we have obtained a cost matrix, we can analyze the shape of the cost matrix to understand the dynamics in the simulations. Furthermore, different cost matrices for different simulations can quantify how the variations of parameters in the simulations affect the dynamics. Lastly, the cost matrix could highlight which physical processes or initial conditions contribute most significantly to the formation of particular structures, such as clusters, filaments or voids. \par 
In this thesis, we will first cover some theoretical foundations of OT and inverse OT in Chapter~\ref{ch:2}, whereafter we will describe the data we will use in more detail in Chapter~\ref{ch:data}. In Chapter~\ref{ch:entropyregiOT}, an algorithm is developed for inverse OT, with comments on the implementation and stability. In Chapter~\ref{ch:results}, the results of the algorithm are discussed for some reference cases, as well as for the FLAMINGO data. 

